% !TeX program = lualatex
\documentclass[11pt]{article}
\usepackage[margin=25mm]{geometry}
\usepackage{fontspec}
\setmainfont{Latin Modern Roman}
\usepackage{amsmath,amssymb,mathtools}
\usepackage{unicode-math}
\setmathfont{Latin Modern Math}
\usepackage{xurl,hyperref}
\hypersetup{colorlinks=true,urlcolor=blue}
\setcounter{secnumdepth}{0}
\setlength{\parindent}{0pt}
\setlength{\parskip}{5pt}
\setlength{\emergencystretch}{3em}
\newcommand{\R}{\mathbb R}
\newcommand{\N}{\mathbb N}
\newcommand{\Z}{\mathbb Z}
\newcommand{\Q}{\mathbb Q}
\newcommand{\C}{\mathbb C}
\newcommand{\D}{\mathbb D}
\newcommand{\F}{\mathbb F}
\newcommand{\sref}[2]{\hyperlink{source:#1:#2}{[S#2]}}
\newcommand{\cataloguescope}[1]{#1}
\begin{document}
\section{Kirby Problem 5.3}\label{statement}

\subsection{Definitions and mathematical statement}
The Hilbert--Smith assertion asks whether every locally compact group acting faithfully and continuously on a connected manifold is a Lie group. The companion question concerns an action of a Cantor group on a Euclidean neighborhood retract: must the subset $A$ of points with trivial stabilizer be a homology-$Z$-subset? This homological negligibility means that $H_*(U,U\setminus A;\mathbb Z)=0$ for every open subset $U$ of the ambient neighborhood retract.
\subsection{Short English statement}
The Hilbert--Smith assertion asks whether every locally compact group acting faithfully and continuously on a connected manifold is a Lie group. The companion question concerns an action of a Cantor group on a Euclidean neighborhood retract: must the subset $A$ of points with trivial stabilizer be a homology-$Z$-subset? This homological negligibility means that $H_*(U,U\setminus A;\mathbb Z)=0$ for every open subset $U$ of the ambient neighborhood retract.
\subsection{Sources}
\begin{itemize}
\item[\hypertarget{source:statement:1}{S1}] ULAM.AI and the UnsolvedMath Contributors, UnsolvedMath: A Curated Collection of Open Mathematics Problems, catalogue entry. Catalogue attribution under CC BY 4.0.
\url{https://huggingface.co/datasets/ulamai/UnsolvedMath}.
\item[\hypertarget{source:statement:2}{S2}] Original problem formulation and mathematical context.
\url{https://math.berkeley.edu/~kirby/problems.pdf}.
\end{itemize}
\label{end:statement}
\end{document}
